\section{Related Work}
\label{sec:related}

\para{Linear representations in language models.}
The linear representation hypothesis~\citep{park2023linear} posits that high-level concepts are encoded as directions in a model's activation space.
\citet{marks2024geometry} demonstrated that truthfulness is linearly represented in LLM residual streams, introducing mass-mean probing (difference-in-means with covariance correction) as a direction-finding method that identifies more causally relevant directions than logistic regression.
\citet{arditi2024refusal} showed that refusal behavior across 13 chat models is mediated by a single direction in the residual stream; erasing this direction disables refusal with minimal capability loss.
\citet{gurnee2023language} found linear representations of spatial and temporal concepts, identifying individual ``space neurons'' and ``time neurons.''
\citet{cunningham2023sparse} showed that sparse autoencoders can decompose activations into interpretable features, resolving superposition.
Our work extends this line of research to a stylistic property---\aistyle---rather than a factual or behavioral one.

\para{Activation steering.}
\citet{turner2023activation} introduced Activation Addition (ActAdd), computing steering vectors from contrastive prompt pairs and adding them during inference to control sentiment and topic.
\citet{rimsky2024steering} scaled this approach with Contrastive Activation Addition (CAA), averaging residual stream differences over many contrastive pairs for more robust steering.
\citet{jorgensen2023improving} showed that mean-centring steering vectors significantly improves effectiveness.
\citet{lee2024mechanistic} found that DPO does not remove toxic capabilities but routes around them, suggesting alignment overlays representations rather than erasing them.
We adopt the contrastive activation methodology from this body of work, adapting it for stylistic rather than behavioral steering.

\para{AI-generated text detection.}
Most approaches to detecting AI-generated text operate on model outputs rather than internal representations.
\citet{guo2023hc3} introduced the \hcthree dataset of paired human and ChatGPT responses and trained surface-level classifiers.
Watermarking methods~\citep{kirchenbauer2023watermark} embed statistical signals in generated text.
Output-probability approaches compare a text's likelihood under different models.
\citet{oneill2025single} is closest to our work: they showed that a single linear direction in an observer model's residual stream can detect hallucinations in another model's outputs.
Unlike all of these approaches, we study whether the \emph{generating} model itself linearly represents the stylistic quality that makes its outputs detectable.

\para{Positioning.}
Our work sits at the intersection of linear representation analysis and AI text detection.
Unlike prior linear representation studies, we target a stylistic property rather than a factual or behavioral one.
Unlike AI detection research, we probe internal representations rather than surface features.
The closest comparison is to the refusal direction work of \citet{arditi2024refusal}: both find directions that mediate a model-level property, but \aistyle turns out to be more complex (multi-dimensional) than refusal (single direction).
